\section{Testresultater}
\label{sec:testresultater}

Testsuiten består av 166 automatiserte tester, og alle passerer i
CI-pipeline. Pipeline kjører \texttt{dotnet test}, men samler ikke inn
coverage. Linje- og branch-dekning ble derfor målt lokalt med Coverlet
(\texttt{--collect:"XPlat Code Coverage"}).

Ved siste lokale måling før innlevering var samlet linjedekning om lag 66\,\%
og branch-dekning om lag 38\,\%. Dekningstallene påvirkes noe av autogenererte
OpenAPI-filer under \texttt{obj/}, som Coverlet teller med som udekket
kode selv om de ikke inneholder forretningslogikk.

Dekningen vurderes som moderat. Den gir god støtte for at sentrale deler
av proof-of-conceptet fungerer, men den er ikke høy nok til å dokumentere
produksjonsmodenhet alene. Gruppen satte heller ikke konkrete dekningsmål
ved prosjektstart. I ettertid vurderes dette som en svakhet, særlig for
kritiske deler som controller-lag, provider-lag, autentisering og
feilhåndtering.

\begin{table}[h]
\centering
\begin{tabular}{llr}
\toprule
\textbf{Komponent} & \textbf{Testtype} & \textbf{Linjedekning} \\
\midrule
AuthController & Integrasjonstest & 100\% \\
KeycloakService & Enhetstest & 100\% \\
VideoProviderFactory & Enhetstest & 100\% \\
AntMediaAuthHandler & Enhetstest & 98\% \\
NHNAuthHandler & Enhetstest & 97,8\% \\
RoomsController & Integrasjonstest & 97\% \\
NHNVideoProvider & Enhetstest & 94,7\% \\
ProviderExtensions & Enhetstest & 94,2\% \\
ApplicationExtensions & Integrasjonstest & 93,7\% \\
GlobalExceptionHandler & Enhetstest & 89,2\% \\
AntMediaProvider & Enhetstest & 78,8\% \\
\midrule
VideoMetrics & Ingen & 0\%$^*$ \\
NHNConference & Ingen & 0\%$^*$ \\
Microsoft.AspNetCore.OpenApi.Generated & Ingen & 0\%$^\dagger$ \\
System.Runtime.CompilerServices & Ingen & 0\%$^\dagger$ \\
\bottomrule
\end{tabular}
\caption{Linjedekning per komponent ved lokal Coverlet-måling.
Samlet linjedekning for løsningen var ca.\ 66\,\%. 
$^*$~Ikke testet direkte. $^\dagger$~Autogenerert kode.}
\label{tab:testresultater}
\end{table}

Tabell~\ref{tab:testresultater} viser høy dekning for flere sentrale
komponenter, blant annet autentisering, provider-factory,
\texttt{RoomsController} og provider-spesifikke auth-handlere. Dette betyr
likevel ikke at alle mulige driftsscenarioer er testet. Dekningstall viser
hvilke linjer som er kjørt, ikke om alle relevante kombinasjoner av input,
nettverksfeil og sikkerhetsscenarioer er undersøkt.

De svakest dekkede områdene er Prometheus-instrumentering,
nettverksnære feil mot eksterne videomotorer, enkelte exception-grener og
mer systematisk negativ input. \texttt{AntMediaProvider} har lavere dekning
enn flere andre komponenter fordi noen grener krever spesifikke
nettverks- eller backend-tilstander som ikke er simulert i enhetstestene.

Manuell verifisering mot HDOs testmiljøer ble gjort med
\texttt{HDO-VideoAPI.http}, utforskende kall og frontend-demoen. Dette
inngår ikke i CI-testsuiten, men ble brukt til å kontrollere komplette
flyter for begge backends, inkludert autentisering, romlivssyklus og
forventede \texttt{501}-responser for NHN-operasjoner som ikke støttes.

